% Logros del laboratorio 4 frente a sus objetivos y trabajo del docente.
% Generado por ../generar_logros.py desde ../logros_datos.py: editar los datos y regenerar.
\documentclass[a4paper,11pt]{article}
% Preámbulo común de los informes de laboratorio (II/2026). Se incluye desde laboratorios/labNN_*/
%   % Preámbulo común de los informes de laboratorio (II/2026). Se incluye desde laboratorios/labNN_*/
%   % Preámbulo común de los informes de laboratorio (II/2026). Se incluye desde laboratorios/labNN_*/
%   \input{../informe_comun}
\input{../../gemelo-digital/manual/estilo_manual}
\input{../../gemelo-digital/manual/tecnico/macros}
\graphicspath{{../../documentacion/planos_isa/pid/}{../../documentacion/planos_isa/electrico/}{../../documentacion/manuales_inversa/}
  {../../database/mockups/}{../../database/documentacion/}{../../database/modelo/}{../../database/interfaz/}
  {../../gemelo-digital/manual/img/}{../../modelo-3d/informe/figuras/}{../../modelo-3d/levantamiento/}{../../../../global/fotos-planta/}{../figuras/}}
\usepackage{pdflscape}
\lstset{showstringspaces=false}
% plano A3 en una página apaisada: \plano[página]{archivo}{leyenda}
\newcommand{\plano}[3][1]{\begin{landscape}\begin{figure}[p]\centering\includegraphics[page=#1, width=\linewidth, height=0.86\textwidth, keepaspectratio]{#2}\caption{#3}\end{figure}\end{landscape}}
\setlength{\headheight}{14pt}
\newcommand{\nlab}{}
\newcommand{\tlab}{}
\lhead{\small Brazo XK-JS1-1 · Informe del laboratorio \nlab}
\rhead{\small UMSA · II/2026}

% \portada{número}{título}{inicio}{entrega}{semanas}
\newcommand{\portada}[5]{%
\renewcommand{\nlab}{#1}%
\begin{titlepage}
\centering
\vspace*{0.6cm}
{\large Universidad Mayor de San Andrés · Facultad de Ingeniería\par}
{\large Carrera de Ingeniería Electrónica\par}
\vspace{0.8cm}
{\color{naranja!80!black}\rule{\linewidth}{1.5pt}\par}
\vspace{0.3cm}
{\Large Informe del Laboratorio Integrado Nº #1\par}
\vspace{0.2cm}
{\huge\bfseries #2\par}
\vspace{0.3cm}
{\color{naranja!80!black}\rule{\linewidth}{1.5pt}\par}
\vspace{0.8cm}
{\large Planta: brazo robótico XK-JS1-1 de 4 grados de libertad\par}
\vspace{0.3cm}
{\large ETN-902 Sistemas de Control II · ETN-1034 Aplicación de Técnicas de Control · ETN-1000 Bases de Datos\par}
\vspace{1cm}
\begin{tabular}{rl}
\textbf{Docente:} & Ing. Jorge Antonio Nava Amador\\
\textbf{Auxiliares:} & Limber Ajnota · Fabián Plata · Franz Choque\\
\textbf{Inicio de trabajo:} & #3\\
\textbf{Entrega:} & #4 (#5)\\
\end{tabular}
\vspace{1cm}

\begin{tabular}{ll}
\multicolumn{2}{l}{\textbf{Escuadra del brazo de 4 GDL}}\\[2pt]
Control (ETN-902 / ETN-1034): & Gabriel Andres Gomez Calla\\
 & Jose Lenin Mamani Choque\\
 & Miguel Alejandro Lima Huarachi\\
 & Román Abner Machicado Arequipa\\[2pt]
Cómputo (ETN-1000): & Miguel Alejandro Lima Huarachi\\
 & Marco Eduardo Chura Rondo\\
\end{tabular}
\vfill
{\small Repositorio: \url{https://github.com/GabAndGom/Brazo-de-4-grados-de-libertad}\par}
{\small La Paz, Bolivia · gestión II/2026\par}
\end{titlepage}
}

% Tabla de entregables (Definition of Done): \begin{dod} ... \end{dod}, filas: entregable & estado & evidencia\\
\newenvironment{dod}{\begin{center}\small\begin{longtable}{>{\raggedright}p{5.2cm} >{\raggedright}p{1.8cm} >{\raggedright\arraybackslash}p{7.6cm}}
\toprule \textbf{Entregable} & \textbf{Estado} & \textbf{Evidencia / observación}\\ \midrule \endhead}
{\bottomrule\end{longtable}\end{center}}

\input{../../gemelo-digital/manual/estilo_manual}
% Macros del manual técnico: ladder, estados, código y rutas de imágenes
\usetikzlibrary{shapes.multipart}
\graphicspath{{../img/}{../../../../../global/fotos-planta/}{../../../modelo-3d/informe/figuras/}{../../../modelo-3d/levantamiento/}{../../../database/modelo/}{../../../documentacion/manuales_inversa/}{../../../documentacion/planos_isa/electrico/}{../../../database/interfaz/}}

\lstdefinestyle{py}{language=Python, keywordstyle=\color{azulg}\bfseries, commentstyle=\color{grisos}\itshape,
  stringstyle=\color{verde}, basicstyle=\ttfamily\scriptsize, numbers=left, numberstyle=\tiny\color{grisos}}
\lstdefinestyle{js}{language=Java, morekeywords={const,let,function,async,await,return,for,of,if,else}, keywordstyle=\color{azulg}\bfseries,
  commentstyle=\color{grisos}\itshape, stringstyle=\color{verde}, basicstyle=\ttfamily\scriptsize, numbers=left, numberstyle=\tiny\color{grisos}}
\lstdefinestyle{awl}{basicstyle=\ttfamily\tiny, commentstyle=\color{grisos}\itshape, morecomment=[l]{//},
  keywordstyle=\color{azulg}\bfseries, morekeywords={LD,LDN,A,AN,O,ON,EU,ED,S,R,MOVB,MOVW,MOVD,PLS,TON,CALL,CRET,LPS,LRD,LPP,ALD,OLD,FILL,LDW,LDB,LDD,AW,AD,AB,OW,OD},
  numbers=left, numberstyle=\tiny\color{grisos}, frame=single, backgroundcolor=\color{gris}}
\lstdefinestyle{sql}{language=SQL, keywordstyle=\color{azulg}\bfseries, commentstyle=\color{grisos}\itshape,
  stringstyle=\color{verde}, basicstyle=\ttfamily\tiny, numbers=left, numberstyle=\tiny\color{grisos},
  morekeywords={REFERENCES,SERIAL,SMALLSERIAL,BIGSERIAL,TIMESTAMPTZ,BOOLEAN,INET,RETURNS,TRIGGER,LANGUAGE,DECLARE,RAISE,EXCEPTION,VIEW,INDEX}}

\tikzset{
  ld/.style={line width=0.7pt},
  pics/na/.style={code={\fill[white](-0.12,-0.26) rectangle (0.12,0.26);
     \draw[thick](-0.1,-0.24)--(-0.1,0.24)(0.1,-0.24)--(0.1,0.24);
     \node[font=\scriptsize\ttfamily, above] at (0,0.22){#1};}},
  pics/nc/.style={code={\fill[white](-0.12,-0.26) rectangle (0.12,0.26);
     \draw[thick](-0.1,-0.24)--(-0.1,0.24)(0.1,-0.24)--(0.1,0.24);\draw(-0.17,-0.24)--(0.17,0.24);
     \node[font=\scriptsize\ttfamily, above] at (0,0.22){#1};}},
  pics/pf/.style={code={\fill[white](-0.12,-0.26) rectangle (0.12,0.26);
     \draw[thick](-0.1,-0.24)--(-0.1,0.24)(0.1,-0.24)--(0.1,0.24);
     \node[font=\tiny\bfseries] at (0,0){P};}},
  pics/cmp/.style n args={3}{code={\fill[white](-0.12,-0.26) rectangle (0.12,0.26);
     \draw[thick](-0.1,-0.24)--(-0.1,0.24)(0.1,-0.24)--(0.1,0.24);
     \node[font=\tiny\ttfamily] at (0,0){#2};
     \node[font=\scriptsize\ttfamily, above] at (0,0.22){#1};
     \node[font=\scriptsize\ttfamily, below] at (0,-0.22){#3};}},
  pics/bob/.style n args={3}{code={\fill[white](-0.2,-0.26) rectangle (0.2,0.26);
     \draw[thick](-0.05,0.24) arc (110:250:0.255);\draw[thick](0.05,0.24) arc (70:-70:0.255);
     \node[font=\tiny\bfseries] at (0,0){#2};
     \node[font=\scriptsize\ttfamily, above] at (0,0.22){#1};
     \node[font=\scriptsize\ttfamily, below] at (0,-0.22){#3};}},
  fb/.style={draw, fill=white, font=\scriptsize\ttfamily, align=left, inner sep=3pt, anchor=west},
  barra/.style={line width=2pt},
  estado/.style={draw, rounded corners=3pt, fill=white, align=center, font=\scriptsize, minimum width=22mm, minimum height=10mm, inner sep=3pt},
  estadon/.style={estado, fill=naranja!14, draw=naranja!70!black},
  trans/.style={font=\scriptsize, align=center, fill=white, inner sep=1pt},
  lugar/.style={circle, draw, thick, minimum size=7mm, fill=white},
  transicion/.style={rectangle, fill=black, minimum width=1.6mm, minimum height=7mm, inner sep=0},
  tabla/.style={rectangle split, rectangle split parts=2, draw, fill=white, font=\tiny\ttfamily, align=left, inner sep=2pt, rounded corners=1pt, text width=33mm},
}
\newcommand{\pendiente}[1]{\par\smallskip\noindent\colorbox{naranja!12}{\parbox{\dimexpr\linewidth-2\fboxsep}{\textbf{\color{naranja!80!black}Pendiente.} #1}}\par\smallskip}
\newcommand{\cumple}{\textcolor{verde}{\textbf{cumplido}}}
\newcommand{\parcial}{\textcolor{naranja}{\textbf{parcial}}}
\newcommand{\falta}{\textcolor{rojo}{\textbf{pendiente}}}

\graphicspath{{../../documentacion/planos_isa/pid/}{../../documentacion/planos_isa/electrico/}{../../documentacion/manuales_inversa/}
  {../../database/mockups/}{../../database/documentacion/}{../../database/modelo/}{../../database/interfaz/}
  {../../gemelo-digital/manual/img/}{../../modelo-3d/informe/figuras/}{../../modelo-3d/levantamiento/}{../../../../global/fotos-planta/}{../figuras/}}
\usepackage{pdflscape}
\lstset{showstringspaces=false}
% plano A3 en una página apaisada: \plano[página]{archivo}{leyenda}
\newcommand{\plano}[3][1]{\begin{landscape}\begin{figure}[p]\centering\includegraphics[page=#1, width=\linewidth, height=0.86\textwidth, keepaspectratio]{#2}\caption{#3}\end{figure}\end{landscape}}
\setlength{\headheight}{14pt}
\newcommand{\nlab}{}
\newcommand{\tlab}{}
\lhead{\small Brazo XK-JS1-1 · Informe del laboratorio \nlab}
\rhead{\small UMSA · II/2026}

% \portada{número}{título}{inicio}{entrega}{semanas}
\newcommand{\portada}[5]{%
\renewcommand{\nlab}{#1}%
\begin{titlepage}
\centering
\vspace*{0.6cm}
{\large Universidad Mayor de San Andrés · Facultad de Ingeniería\par}
{\large Carrera de Ingeniería Electrónica\par}
\vspace{0.8cm}
{\color{naranja!80!black}\rule{\linewidth}{1.5pt}\par}
\vspace{0.3cm}
{\Large Informe del Laboratorio Integrado Nº #1\par}
\vspace{0.2cm}
{\huge\bfseries #2\par}
\vspace{0.3cm}
{\color{naranja!80!black}\rule{\linewidth}{1.5pt}\par}
\vspace{0.8cm}
{\large Planta: brazo robótico XK-JS1-1 de 4 grados de libertad\par}
\vspace{0.3cm}
{\large ETN-902 Sistemas de Control II · ETN-1034 Aplicación de Técnicas de Control · ETN-1000 Bases de Datos\par}
\vspace{1cm}
\begin{tabular}{rl}
\textbf{Docente:} & Ing. Jorge Antonio Nava Amador\\
\textbf{Auxiliares:} & Limber Ajnota · Fabián Plata · Franz Choque\\
\textbf{Inicio de trabajo:} & #3\\
\textbf{Entrega:} & #4 (#5)\\
\end{tabular}
\vspace{1cm}

\begin{tabular}{ll}
\multicolumn{2}{l}{\textbf{Escuadra del brazo de 4 GDL}}\\[2pt]
Control (ETN-902 / ETN-1034): & Gabriel Andres Gomez Calla\\
 & Jose Lenin Mamani Choque\\
 & Miguel Alejandro Lima Huarachi\\
 & Román Abner Machicado Arequipa\\[2pt]
Cómputo (ETN-1000): & Miguel Alejandro Lima Huarachi\\
 & Marco Eduardo Chura Rondo\\
\end{tabular}
\vfill
{\small Repositorio: \url{https://github.com/GabAndGom/Brazo-de-4-grados-de-libertad}\par}
{\small La Paz, Bolivia · gestión II/2026\par}
\end{titlepage}
}

% Tabla de entregables (Definition of Done): \begin{dod} ... \end{dod}, filas: entregable & estado & evidencia\\
\newenvironment{dod}{\begin{center}\small\begin{longtable}{>{\raggedright}p{5.2cm} >{\raggedright}p{1.8cm} >{\raggedright\arraybackslash}p{7.6cm}}
\toprule \textbf{Entregable} & \textbf{Estado} & \textbf{Evidencia / observación}\\ \midrule \endhead}
{\bottomrule\end{longtable}\end{center}}

\input{../../gemelo-digital/manual/estilo_manual}
% Macros del manual técnico: ladder, estados, código y rutas de imágenes
\usetikzlibrary{shapes.multipart}
\graphicspath{{../img/}{../../../../../global/fotos-planta/}{../../../modelo-3d/informe/figuras/}{../../../modelo-3d/levantamiento/}{../../../database/modelo/}{../../../documentacion/manuales_inversa/}{../../../documentacion/planos_isa/electrico/}{../../../database/interfaz/}}

\lstdefinestyle{py}{language=Python, keywordstyle=\color{azulg}\bfseries, commentstyle=\color{grisos}\itshape,
  stringstyle=\color{verde}, basicstyle=\ttfamily\scriptsize, numbers=left, numberstyle=\tiny\color{grisos}}
\lstdefinestyle{js}{language=Java, morekeywords={const,let,function,async,await,return,for,of,if,else}, keywordstyle=\color{azulg}\bfseries,
  commentstyle=\color{grisos}\itshape, stringstyle=\color{verde}, basicstyle=\ttfamily\scriptsize, numbers=left, numberstyle=\tiny\color{grisos}}
\lstdefinestyle{awl}{basicstyle=\ttfamily\tiny, commentstyle=\color{grisos}\itshape, morecomment=[l]{//},
  keywordstyle=\color{azulg}\bfseries, morekeywords={LD,LDN,A,AN,O,ON,EU,ED,S,R,MOVB,MOVW,MOVD,PLS,TON,CALL,CRET,LPS,LRD,LPP,ALD,OLD,FILL,LDW,LDB,LDD,AW,AD,AB,OW,OD},
  numbers=left, numberstyle=\tiny\color{grisos}, frame=single, backgroundcolor=\color{gris}}
\lstdefinestyle{sql}{language=SQL, keywordstyle=\color{azulg}\bfseries, commentstyle=\color{grisos}\itshape,
  stringstyle=\color{verde}, basicstyle=\ttfamily\tiny, numbers=left, numberstyle=\tiny\color{grisos},
  morekeywords={REFERENCES,SERIAL,SMALLSERIAL,BIGSERIAL,TIMESTAMPTZ,BOOLEAN,INET,RETURNS,TRIGGER,LANGUAGE,DECLARE,RAISE,EXCEPTION,VIEW,INDEX}}

\tikzset{
  ld/.style={line width=0.7pt},
  pics/na/.style={code={\fill[white](-0.12,-0.26) rectangle (0.12,0.26);
     \draw[thick](-0.1,-0.24)--(-0.1,0.24)(0.1,-0.24)--(0.1,0.24);
     \node[font=\scriptsize\ttfamily, above] at (0,0.22){#1};}},
  pics/nc/.style={code={\fill[white](-0.12,-0.26) rectangle (0.12,0.26);
     \draw[thick](-0.1,-0.24)--(-0.1,0.24)(0.1,-0.24)--(0.1,0.24);\draw(-0.17,-0.24)--(0.17,0.24);
     \node[font=\scriptsize\ttfamily, above] at (0,0.22){#1};}},
  pics/pf/.style={code={\fill[white](-0.12,-0.26) rectangle (0.12,0.26);
     \draw[thick](-0.1,-0.24)--(-0.1,0.24)(0.1,-0.24)--(0.1,0.24);
     \node[font=\tiny\bfseries] at (0,0){P};}},
  pics/cmp/.style n args={3}{code={\fill[white](-0.12,-0.26) rectangle (0.12,0.26);
     \draw[thick](-0.1,-0.24)--(-0.1,0.24)(0.1,-0.24)--(0.1,0.24);
     \node[font=\tiny\ttfamily] at (0,0){#2};
     \node[font=\scriptsize\ttfamily, above] at (0,0.22){#1};
     \node[font=\scriptsize\ttfamily, below] at (0,-0.22){#3};}},
  pics/bob/.style n args={3}{code={\fill[white](-0.2,-0.26) rectangle (0.2,0.26);
     \draw[thick](-0.05,0.24) arc (110:250:0.255);\draw[thick](0.05,0.24) arc (70:-70:0.255);
     \node[font=\tiny\bfseries] at (0,0){#2};
     \node[font=\scriptsize\ttfamily, above] at (0,0.22){#1};
     \node[font=\scriptsize\ttfamily, below] at (0,-0.22){#3};}},
  fb/.style={draw, fill=white, font=\scriptsize\ttfamily, align=left, inner sep=3pt, anchor=west},
  barra/.style={line width=2pt},
  estado/.style={draw, rounded corners=3pt, fill=white, align=center, font=\scriptsize, minimum width=22mm, minimum height=10mm, inner sep=3pt},
  estadon/.style={estado, fill=naranja!14, draw=naranja!70!black},
  trans/.style={font=\scriptsize, align=center, fill=white, inner sep=1pt},
  lugar/.style={circle, draw, thick, minimum size=7mm, fill=white},
  transicion/.style={rectangle, fill=black, minimum width=1.6mm, minimum height=7mm, inner sep=0},
  tabla/.style={rectangle split, rectangle split parts=2, draw, fill=white, font=\tiny\ttfamily, align=left, inner sep=2pt, rounded corners=1pt, text width=33mm},
}
\newcommand{\pendiente}[1]{\par\smallskip\noindent\colorbox{naranja!12}{\parbox{\dimexpr\linewidth-2\fboxsep}{\textbf{\color{naranja!80!black}Pendiente.} #1}}\par\smallskip}
\newcommand{\cumple}{\textcolor{verde}{\textbf{cumplido}}}
\newcommand{\parcial}{\textcolor{naranja}{\textbf{parcial}}}
\newcommand{\falta}{\textcolor{rojo}{\textbf{pendiente}}}

\graphicspath{{../../documentacion/planos_isa/pid/}{../../documentacion/planos_isa/electrico/}{../../documentacion/manuales_inversa/}
  {../../database/mockups/}{../../database/documentacion/}{../../database/modelo/}{../../database/interfaz/}
  {../../gemelo-digital/manual/img/}{../../modelo-3d/informe/figuras/}{../../modelo-3d/levantamiento/}{../../../../global/fotos-planta/}{../figuras/}}
\usepackage{pdflscape}
\lstset{showstringspaces=false}
% plano A3 en una página apaisada: \plano[página]{archivo}{leyenda}
\newcommand{\plano}[3][1]{\begin{landscape}\begin{figure}[p]\centering\includegraphics[page=#1, width=\linewidth, height=0.86\textwidth, keepaspectratio]{#2}\caption{#3}\end{figure}\end{landscape}}
\setlength{\headheight}{14pt}
\newcommand{\nlab}{}
\newcommand{\tlab}{}
\lhead{\small Brazo XK-JS1-1 · Informe del laboratorio \nlab}
\rhead{\small UMSA · II/2026}

% \portada{número}{título}{inicio}{entrega}{semanas}
\newcommand{\portada}[5]{%
\renewcommand{\nlab}{#1}%
\begin{titlepage}
\centering
\vspace*{0.6cm}
{\large Universidad Mayor de San Andrés · Facultad de Ingeniería\par}
{\large Carrera de Ingeniería Electrónica\par}
\vspace{0.8cm}
{\color{naranja!80!black}\rule{\linewidth}{1.5pt}\par}
\vspace{0.3cm}
{\Large Informe del Laboratorio Integrado Nº #1\par}
\vspace{0.2cm}
{\huge\bfseries #2\par}
\vspace{0.3cm}
{\color{naranja!80!black}\rule{\linewidth}{1.5pt}\par}
\vspace{0.8cm}
{\large Planta: brazo robótico XK-JS1-1 de 4 grados de libertad\par}
\vspace{0.3cm}
{\large ETN-902 Sistemas de Control II · ETN-1034 Aplicación de Técnicas de Control · ETN-1000 Bases de Datos\par}
\vspace{1cm}
\begin{tabular}{rl}
\textbf{Docente:} & Ing. Jorge Antonio Nava Amador\\
\textbf{Auxiliares:} & Limber Ajnota · Fabián Plata · Franz Choque\\
\textbf{Inicio de trabajo:} & #3\\
\textbf{Entrega:} & #4 (#5)\\
\end{tabular}
\vspace{1cm}

\begin{tabular}{ll}
\multicolumn{2}{l}{\textbf{Escuadra del brazo de 4 GDL}}\\[2pt]
Control (ETN-902 / ETN-1034): & Gabriel Andres Gomez Calla\\
 & Jose Lenin Mamani Choque\\
 & Miguel Alejandro Lima Huarachi\\
 & Román Abner Machicado Arequipa\\[2pt]
Cómputo (ETN-1000): & Miguel Alejandro Lima Huarachi\\
 & Marco Eduardo Chura Rondo\\
\end{tabular}
\vfill
{\small Repositorio: \url{https://github.com/GabAndGom/Brazo-de-4-grados-de-libertad}\par}
{\small La Paz, Bolivia · gestión II/2026\par}
\end{titlepage}
}

% Tabla de entregables (Definition of Done): \begin{dod} ... \end{dod}, filas: entregable & estado & evidencia\\
\newenvironment{dod}{\begin{center}\small\begin{longtable}{>{\raggedright}p{5.2cm} >{\raggedright}p{1.8cm} >{\raggedright\arraybackslash}p{7.6cm}}
\toprule \textbf{Entregable} & \textbf{Estado} & \textbf{Evidencia / observación}\\ \midrule \endhead}
{\bottomrule\end{longtable}\end{center}}

\renewcommand{\tipodoc}{Logros del Laboratorio Integrado}
\renewcommand{\tipocab}{Logros del laboratorio}
\begin{document}
\portada{4}{Identificación de la planta y estructura física de la base de datos}{martes 15 de septiembre de 2026}{martes 29 de septiembre de 2026}{semanas 7 y 8}
\tableofcontents
\newpage
\nota{\textbf{Estado al 18/11/2026 (actualizado por la noche).} Laboratorio entregado; la base se ejecutó en PostgreSQL el 12/11/2026. Fase 2: modelado dinámico, caja negra y backend. Este documento explica, objetivo por objetivo, qué pide la guía, qué logró la escuadra, qué es cada cosa y dónde está la evidencia; después explica el trabajo del docente: qué debía revisar en cada hito y qué tiene preparado la escuadra para esa revisión.}
\section{Resumen}
\begin{center}\small\begin{tabularx}{\linewidth}{c >{\raggedright}p{2.4cm} >{\raggedright}X c c}
\toprule \textbf{Nº} & \textbf{Mención} & \textbf{Objetivo} & \textbf{Estado} & \textbf{\%}\\ \midrule
1 & ETN-902 / ETN-1034 & Identificación experimental por caja negra & \parcial & 40\%\\
2 & ETN-902 / ETN-1034 & Modelado analítico multivariable & \parcial & 80\%\\
3 & ETN-1000 & Estructuración física de la base de datos & \parcial & 85\%\\
4 & ETN-1000 & Diccionarios y restricciones & \cumple & 100\%\\
\bottomrule\end{tabularx}\end{center}
El porcentaje mide cuánto del objetivo está hecho y comprobado en la planta o en el repositorio; no cuenta lo que solo está planificado.
\section{Logros respecto de los objetivos}
\subsection{Mención Control Automático (ETN-902 / ETN-1034)}
\subsubsection{Identificación experimental por caja negra}
\paragraph{Qué pide la guía.} Aplicar pruebas escalón a la planta y aproximar su función de transferencia.
\paragraph{Qué se logró.}
\begin{itemize}
\item Procedimiento de la prueba escrito: homing, escalón de frecuencia del tren de pulsos (1000 a 4000 Hz), registro en CSV del tiempo, la frecuencia, los pulsos y el instante en que se activa el final de carrera.
\item Explicación de por qué la curva de reacción no se puede medir hoy: el brazo no tiene sensor de posición; la posición mostrada es la cuenta de pulsos, es decir, la propia entrada integrada.
\item Prueba automatizada \codigo{prueba\_escalon.py}: toma el mando, aplica el escalón de frecuencia, guarda el CSV crudo y calcula K y el tiempo muerto t\_d a partir de las banderas del PLC. Probada con el PLC simulado (4 ensayos, K/K\_nom = 0,998 a 1,000); modelo identificado K$\cdot$e\^{}(-td$\cdot$s)/s.
\end{itemize}
\paragraph{Qué es y por qué importa.} La identificación por caja negra excita la planta con un escalón y mide cómo responde para estimar K, $\tau$ y el tiempo muerto de G(s) = K$\cdot$e\^{}(-td$\cdot$s)/($\tau$s+1), sin conocer su física. Necesita medir la salida y(t) de forma independiente.
\paragraph{Evidencia.}
\begin{itemize}\small
\item \codigo{laboratorios/lab04\_identificacion\_y\_ddl/informe\_lab04.pdf}
\item \codigo{laboratorios/lab04\_identificacion\_y\_ddl/prueba\_escalon.py}
\end{itemize}
\paragraph{Lo que falta.} Correr prueba\_escalon.py con la planta real; para la curva completa ($\tau$) hace falta un encoder o una regla lineal.
\paragraph{Estado:} \parcial~(40\%).
\subsubsection{Modelado analítico multivariable}
\paragraph{Qué pide la guía.} Formular las ecuaciones del sistema y representarlo en espacio de estados $\dot{x}$ = Ax + Bu.
\paragraph{Qué se logró.}
\begin{itemize}
\item Leyes físicas: cada eje lo mueve un motor paso a paso; mientras no pierda pasos, la velocidad es proporcional a la frecuencia de pulsos (190,3 pulsos/° medidos en la base; 200 pulsos/mm en los carros, SUPUESTO por el paso del husillo).
\item Estado x = [$\theta$, sX, sY], entrada u = [f0, f1], A = 0 y B(K1, K2) de 3$\times$2 que depende de los relés (sistema conmutado).
\item Salida C = jacobiano de la cinemática del brazo en el centro del espacio de trabajo, calculado con las cotas del modelo de Inventor del equipo; D = 0.
\end{itemize}
\paragraph{Qué es y por qué importa.} El espacio de estados describe el sistema con un vector x de magnitudes internas y matrices que dicen cómo cambian. Aquí A = 0 significa que sin pulsos ningún eje se mueve (el motor retiene), y B dice cuánto avanza cada eje por hercio; cada eje es un integrador.
\paragraph{Evidencia.}
\begin{itemize}\small
\item \codigo{laboratorios/lab04\_identificacion\_y\_ddl/informe\_lab04.pdf}
\item \codigo{laboratorios/lab05\_gemelo\_y\_triggers/espacio\_estados\_brazo.py}
\end{itemize}
\paragraph{Lo que falta.} Modelo dinámico con inercia y fricción, que requiere medir esos parámetros.
\paragraph{Estado:} \parcial~(80\%).
\subsection{Mención Sistemas Computacionales (ETN-1000)}
\subsubsection{Estructuración física de la base de datos}
\paragraph{Qué pide la guía.} Llevar el modelo lógico a PostgreSQL con un script DDL.
\paragraph{Qué se logró.}
\begin{itemize}
\item Script \codigo{brazo\_4gdl.sql} que crea en una transacción las 15 tablas, 2 índices por marca de tiempo, el trigger de seguridad, 2 vistas y los datos fijos de la planta.
\item Ejecutado sin errores en PostgreSQL 16 en la PC de desarrollo (base \codigo{brazo4gdl}), con registro en \codigo{registro\_ejecucion.txt}.
\end{itemize}
\paragraph{Qué es y por qué importa.} El DDL (Data Definition Language) son las sentencias CREATE TABLE que convierten el diagrama en una base real. Un único script permite crear la misma base en cualquier PC con un comando: \codigo{psql -d brazo4gdl -f brazo\_4gdl.sql}.
\paragraph{Evidencia.}
\begin{itemize}\small
\item \codigo{database/scripts/brazo\_4gdl.sql}
\item \codigo{database/scripts/registro\_ejecucion.txt}
\end{itemize}
\paragraph{Lo que falta.} Ejecutarlo en la PC del laboratorio asignada a la planta.
\paragraph{Estado:} \parcial~(85\%).
\subsubsection{Diccionarios y restricciones}
\paragraph{Qué pide la guía.} Definir tipos de datos, PK, FK y NOT NULL.
\paragraph{Qué se logró.}
\begin{itemize}
\item Tipos elegidos por uso: DOUBLE PRECISION para valores medidos, NUMERIC(10,4) para calibraciones, TIMESTAMPTZ para marcas de tiempo, BOOLEAN para estados, INET para la IP de cada sesión.
\item 18 claves foráneas con ON DELETE RESTRICT y ON UPDATE CASCADE, NOT NULL en todo dato obligatorio, CHECK de rangos y modos, UNIQUE en los tags.
\item Diccionario de datos generado automáticamente del script (\codigo{generar\_diccionario.py}).
\end{itemize}
\paragraph{Qué es y por qué importa.} Las restricciones hacen que la propia base rechace datos imposibles: un tag repetido, una muestra sin variable o borrar un sensor que tiene historial. ON DELETE RESTRICT impide perder historial; ON UPDATE CASCADE propaga un cambio de clave.
\paragraph{Evidencia.}
\begin{itemize}\small
\item \codigo{database/documentacion/base\_datos\_brazo.pdf}
\item \codigo{database/documentacion/generar\_diccionario.py}
\end{itemize}
\paragraph{Lo que falta.} nada para este objetivo.
\paragraph{Estado:} \cumple~(100\%).
\section{Trabajo del docente}
La guía asigna al docente el papel de auditor: en cada hito revisa el trabajo de la escuadra y decide si cumple antes de pasar a la etapa siguiente. Para cada tarea se explica qué debía revisar, cuándo, y qué tiene la escuadra para presentarle. El porcentaje es el del Excel \codigo{evaluacion\_tareas\_docente.xlsx}.
\subsection{Auditoría de modelos (matrices A, B, C, D)}
\paragraph{Lo que dice la guía.} Evaluar de forma personalizada la consistencia analítica de las matrices A, B, C, D. \textit{Cuándo:} martes 29 de septiembre de 2026.
\paragraph{Qué debía revisar el docente (hitos de la revisión).}
\begin{itemize}
\item[$\square$] Que cada componente del vector de estado tenga significado físico y unidad.
\item[$\square$] Que A y B salgan de leyes físicas de la planta y no de un ejemplo genérico.
\item[$\square$] Que el modelo sea coherente con lo que se observa (30 \% de la defensa).
\end{itemize}
\paragraph{Qué tiene la escuadra para esa revisión.}
\begin{itemize}
\item A = 0, B(K1, K2) de 3$\times$2, C = jacobiano en el centro, D = 0, con su deducción en el informe 4.
\item Rangos de controlabilidad y observabilidad calculados con \codigo{espacio\_estados\_brazo.py} / \codigo{.m}.
\end{itemize}
\paragraph{Lo que falta para que el docente lo apruebe.} Modelo dinámico e identificación experimental.
\paragraph{Estado:} \parcial~(60\%).
\subsection{Verificación del backend}
\paragraph{Lo que dice la guía.} Comprobar la estructura física y la normalización del esquema inyectado en las PCs del laboratorio. \textit{Cuándo:} martes 29 de septiembre de 2026.
\paragraph{Qué debía revisar el docente (hitos de la revisión).}
\begin{itemize}
\item[$\square$] Que el script corra sin errores en la PC de la planta.
\item[$\square$] Que las FK tengan RESTRICT/CASCADE y los tipos sean coherentes (40 \% de la defensa).
\item[$\square$] Que el registro de la ejecución esté en \codigo{/database/scripts/}.
\end{itemize}
\paragraph{Qué tiene la escuadra para esa revisión.}
\begin{itemize}
\item \codigo{brazo\_4gdl.sql} ejecutado sin errores en PostgreSQL 16 (15 tablas, 18 FK, 2 índices, trigger, 2 vistas) en la PC de desarrollo.
\end{itemize}
\paragraph{Lo que falta para que el docente lo apruebe.} Ejecutarlo en la PC del laboratorio.
\paragraph{Estado:} \parcial~(75\%).
\subsection{Cápsulas MATLAB/Simulink (tarea de la fase 2)}
\paragraph{Lo que dice la guía.} Proporcionar cápsulas de problemas para identificación por caja negra. \textit{Cuándo:} fase 2 (semanas 7 a 13).
\paragraph{Qué debía revisar el docente (hitos de la revisión).}
\begin{itemize}
\item[$\square$] Que cada escuadra adapte la cápsula a su planta y la ejecute con datos propios.
\end{itemize}
\paragraph{Qué tiene la escuadra para esa revisión.}
\begin{itemize}
\item Modelo analítico calculado en \codigo{espacio\_estados\_brazo.py} y \codigo{.m}; procedimiento de la prueba escalón en el informe 4.
\end{itemize}
\paragraph{Estado:} \parcial~(75\%).
\subsection{Revisar y firmar los esquemas de base de datos (tarea de la fase 2)}
\paragraph{Lo que dice la guía.} Revisar y firmar los esquemas propuestos, validando su modularidad. \textit{Cuándo:} fase 2 (semanas 7 a 13).
\paragraph{Qué debía revisar el docente (hitos de la revisión).}
\begin{itemize}
\item[$\square$] Que el esquema esté separado en bloques independientes (catálogo, operación, telemetría, alarmas, usuarios) y que se pueda extender a otras plantas.
\end{itemize}
\paragraph{Qué tiene la escuadra para esa revisión.}
\begin{itemize}
\item Esquema 3FN de 15 tablas con su documento de 25 páginas, ejecutado con 7/7 pruebas correctas.
\end{itemize}
\paragraph{Lo que falta para que el docente lo apruebe.} Firma del docente.
\paragraph{Estado:} \parcial~(90\%).
\section{Conclusión}
Cumplidos: Diccionarios y restricciones. Con avance parcial: Identificación experimental por caja negra, Modelado analítico multivariable, Estructuración física de la base de datos.
\par\smallskip Pendientes principales:
\begin{enumerate}
\item Correr prueba\_escalon.py con la planta real; para la curva completa ($\tau$) hace falta un encoder o una regla lineal.
\item Modelo dinámico con inercia y fricción, que requiere medir esos parámetros.
\item Ejecutarlo en la PC del laboratorio asignada a la planta.
\end{enumerate}
\par\smallskip\noindent Repositorio: \url{https://github.com/GabAndGom/Brazo-de-4-grados-de-libertad}
\end{document}
